% Generated from validated Special post-draft synthesis. Do not hand-edit.
\section*{Monthly Signals}
\addcontentsline{toc}{section}{Monthly Signals}
\smallskip
\noindent \textbf{Frontier Models \& Developer Platform} model更新はtool、computer use、multimodal、API lifecycleという利用surfaceの変化と同時に進んだ。 \hfill{\footnotesize p.~\pageref{pkg:frontier-models-platform-march}}
\par
\smallskip
\noindent \textbf{Inference, Serving \& Systems} cache、routing、communication、RPC、tool-aware servingが分化し、modelの外側の実行基盤が独立した競争領域になった。 \hfill{\footnotesize p.~\pageref{pkg:inference-serving-systems-march}}
\par
\smallskip
\noindent \textbf{Safety, Security \& Control} prompt injection対策はinstruction hierarchy、tool privilege、configuration、monitoringまで含む多層のsystem designへ広がった。 \hfill{\footnotesize p.~\pageref{pkg:safety-security-control-march}}
\par
\smallskip
\noindent \textbf{Multimodal \& Generation} generationとunderstandingの統合、grounding、on-device deploymentが並行して進み、multimodalは独立した主要テーマとして厚みを持った。 \hfill{\footnotesize p.~\pageref{pkg:multimodal-generation-march}}
\par
\smallskip
\noindent \textbf{Agents, Runtime \& Evaluation} Agentは最大テーマではない一方、execution environment、long-horizon reliability、memory、verificationと、その評価方法をめぐる論点が具体化した。 \hfill{\footnotesize p.~\pageref{pkg:agents-runtime-long-horizon-march} / p.~\pageref{pkg:paper-watch-march}}
\par
\medskip
\begin{claimboundary}[Retrospective scope]
本号は2026年3月を後日確認可能になった一次情報も用いて再構成するRetrospective Specialである。Coverage Periodと制作時点を同一視せず、vendor / project / author claimの境界は各記事で明示する。
\end{claimboundary}
\medskip
\tableofcontents
